\documentclass[12pt]{article}
\usepackage[utf8]{inputenc}
\usepackage{amsmath}
\usepackage{graphicx}
\usepackage{natbib}
\usepackage{hyperref}
\usepackage{lipsum} % For dummy text, you can remove this in your final document

\title{Biological Mechanisms of Aging and Longevity: Genetic, Epigenetic, and Environmental Influences}
\author{Your Name}
\date{\today}

\begin{document}

\maketitle

\begin{abstract}
Understanding the biological mechanisms underlying aging and longevity is crucial for addressing age-related diseases and promoting healthy aging. This paper reviews genetic, epigenetic, and environmental factors influencing lifespan and healthspan, discusses current interventions, and explores future research directions in the field of aging biology.
\end{abstract}

\section{Introduction}
Aging is a complex biological process characterized by the progressive decline in physiological function and resilience over time. This paper explores the mechanisms of aging at the molecular and cellular levels, examines genetic, epigenetic, and environmental influences on longevity, and discusses the implications for aging-related diseases and interventions.

\section{Mechanisms of Aging}
\subsection{Genetic Factors}
Genetic variants and mutations influence aging processes, affecting pathways such as DNA repair mechanisms, telomere maintenance, and mitochondrial function \citep{kenyon2010mechanisms, lopez-otin2013hallmarks}.

\subsection{Epigenetic Modifications}
Epigenetic changes, including DNA methylation, histone modifications, and non-coding RNA regulation, alter gene expression patterns with age and impact cellular senescence and longevity \citep{booth2016epigenetic, maegawa2017epigenetic}.

\subsection{Cellular Senescence}
Accumulation of senescent cells contributes to tissue dysfunction and age-related diseases through inflammation and impaired regeneration mechanisms \citep{he2019cellular, demaria2014cellular}.

\section{Influencing Factors}
\subsection{Environmental Influences}
External factors such as diet, exercise, exposure to toxins, and socio-economic status influence aging processes and health outcomes across the lifespan \citep{kirkland2017aging, franceschi2018inflammaging}.

\subsection{Lifestyle and Behavioral Factors}
Health behaviors, including smoking, alcohol consumption, stress management, and social interactions, impact aging trajectories and longevity outcomes \citep{willcox2007healthy, stewart2018stress}.

\section{Interventions}
\subsection{Nutritional Interventions}
Caloric restriction, dietary supplementation with antioxidants and anti-inflammatory compounds, and personalized nutrition approaches mitigate aging-related cellular damage and promote longevity \citep{mattson2014nutritional, fontana2014longevity}.

\subsection{Pharmacological Interventions}
Targeted drug therapies, including senolytics and geroprotectors, aim to delay aging processes and reduce age-related diseases by targeting specific molecular pathways \citep{kirkland2015senolytics, kennedy2016geroscience}.

\section{Future Research}
Future research directions include elucidating novel aging pathways, integrating multi-omics approaches for personalized aging interventions, and exploring ethical and societal implications of extending healthspan and lifespan.

\section*{References}
\bibliographystyle{plainnat}
\bibliography{prompt9_35}

\end{document}
